\documentclass[a3paper, landscape, 12pt]{scrartcl}

% Exakte Ränder: 1.2cm rundum, keine Seitenzahlen
\usepackage[a3paper, landscape, margin=1.2cm]{geometry}
\usepackage[ngerman]{babel}
\usepackage{libertinus}
\usepackage{xcolor}
\usepackage{tikz}

\definecolor{primaryBlue}{RGB}{20, 50, 100}
\definecolor{lightBlue}{RGB}{245, 248, 253}
\definecolor{accentBlue}{RGB}{37, 99, 235}

\pagestyle{empty}
\setlength{\parindent}{0pt}

\begin{document}%
\noindent%
\begin{tikzpicture}[x=1cm, y=1cm]

    % =========================================================================
    % 1. GEOMETRIE (EXAKT NACH DER SKIZZE DES NUTZERS)
    % Bounding Box: 39.0cm x 26.5cm (passt perfekt auf 1 Seite A3)
    % =========================================================================

    % Eckpunkte Außenrahmen
    \coordinate (BL) at (0, 0);
    \coordinate (BR) at (39.0, 0);
    \coordinate (TR) at (39.0, 26.5);
    \coordinate (TL) at (0, 26.5);

    % Eckpunkte Innenkasten (zentriert: Breite 20.0cm, Höhe 12.0cm)
    % x von 9.5 bis 29.5 (Mitte bei 19.5)
    % y von 7.2 bis 19.2 (Mitte bei 13.2)
    \coordinate (IBL) at (9.5, 7.2);
    \coordinate (IBR) at (29.5, 7.2);
    \coordinate (ITR) at (29.5, 19.2);
    \coordinate (ITL) at (9.5, 19.2);

    % Mittelpunkte der unteren Kanten
    \coordinate (IBC) at (19.5, 7.2);
    \coordinate (BC)  at (19.5, 0);

    % Hintergrund für das Mittelfeld
    \fill[color=lightBlue!50] (IBL) rectangle (ITR);

    % --- LINIEN (EXAKT WIE AUF DER VORLAGE) ---
    % 1. Außenrahmen
    \draw[line width=1.8pt, color=primaryBlue] (BL) rectangle (TR);

    % 2. Innenkasten
    \draw[line width=1.8pt, color=primaryBlue] (IBL) rectangle (ITR);

    % 3. Diagonale Oben-Links
    \draw[line width=1.4pt, color=primaryBlue] (TL) -- (ITL);

    % 4. Diagonale Oben-Rechts
    \draw[line width=1.4pt, color=primaryBlue] (TR) -- (ITR);

    % 5. Senkrechte Unten-Mitte
    \draw[line width=1.4pt, color=primaryBlue] (IBC) -- (BC);


    % =========================================================================
    % 2. INHALT MITTELFELD: GEMEINSAME STRUKTURLEGETECHNIK
    % =========================================================================
    \node[anchor=north, text width=19.0cm, align=center, font=\sffamily] at (19.5, 18.7) {
        {\Large\bfseries\color{primaryBlue} Mittelfeld: Gemeinsame Strukturlegetechnik}\\[0.15cm]
        {\small\color{primaryBlue!85}\textbf{Anleitung:} Einigt euch auf die \textbf{5 Kernbegriffe}. Legt und ordnet sie als \textbf{Ober- und Unterbegriffe} bzw. in ihr Wirkungsgefüge und verbindet sie mit beschrifteten Pfeilen ($\rightarrow$, $\leftrightarrow$).}
    };

    % Zarte Orientierungs-Kästchen für die 5 Stationen (lassen viel Platz zum Legen)
    \node[anchor=center, text width=19.0cm, align=center, font=\footnotesize\sffamily\color{gray!70}] at (19.5, 12.5) {
        \begin{tikzpicture}[x=3.7cm, y=1.2cm]
            \node[draw=primaryBlue!35, rounded corners, fill=white, inner sep=5pt, text width=3.2cm, align=center] (b1) at (0, 1.5) {\textbf{\color{primaryBlue}1. Biologische Natur}\\ \tiny Triebbegabung / Homo homini lupus};
            \node[draw=primaryBlue!35, rounded corners, fill=white, inner sep=5pt, text width=3.2cm, align=center] (b2) at (2.4, 1.5) {\textbf{\color{primaryBlue}2. Bedrohte Kultur}\\ \tiny Zerfall der Gesellschaft};
            \node[draw=primaryBlue!35, rounded corners, fill=white, inner sep=5pt, text width=3.4cm, align=center] (b3) at (4.8, 1.5) {\textbf{\color{primaryBlue}3. Kulturelle Zähmung}\\ \tiny Triebverzicht / Liebesgebot};
            \node[draw=primaryBlue!35, rounded corners, fill=white, inner sep=5pt, text width=3.4cm, align=center] (b4) at (1.2, -1.5) {\textbf{\color{primaryBlue}4. Psychischer Wandel}\\ \tiny Introjektion $\rightarrow$ Über-Ich};
            \node[draw=primaryBlue!35, rounded corners, fill=white, inner sep=5pt, text width=3.6cm, align=center] (b5) at (3.6, -1.5) {\textbf{\color{primaryBlue}5. Seelische Kostenstelle}\\ \tiny Schuldbewusstsein / Unbehagen};

            \draw[->, >=latex, line width=0.8pt, color=primaryBlue!35, dashed] (b1) -- node[above, font=\tiny] {erzeugt} (b2);
            \draw[->, >=latex, line width=0.8pt, color=primaryBlue!35, dashed] (b2) -- node[above, font=\tiny] {erzwingt} (b3);
            \draw[->, >=latex, line width=0.8pt, color=primaryBlue!35, dashed] (b3) -- node[right, font=\tiny] {führt zu} (b4);
            \draw[->, >=latex, line width=0.8pt, color=primaryBlue!35, dashed] (b4) -- node[above, font=\tiny] {erzeugt} (b5);
        \end{tikzpicture}
    };


    % =========================================================================
    % 3. BESCHRIFTUNG DER 3 AUSSENFELDER (FÜR 3ER-GRUPPE)
    % =========================================================================

    % --- FELD 1: OBEN (Trapez) ---
    \node[anchor=north, text width=26cm, align=center, font=\sffamily] at (19.5, 25.7) {
        \textbf{\color{primaryBlue}\large Außenfeld 1 (Schüler/in A)} \qquad Name: \rule{5cm}{0.4pt}\\[0.15cm]
        \small\textbf{1.} Lies den Textauszug und notiere für dich \textbf{5--8 zentrale Begriffe}. \quad
        \textbf{2.} Lies die Felder der anderen und markiere: \textbf{!} = tragend, \textbf{?} = unklar.
    };

    % --- FELD 2: UNTEN LINKS ---
    \node[anchor=north west, text width=17.0cm, align=left, font=\sffamily] at (1.0, 6.4) {
        \textbf{\color{primaryBlue}\large Außenfeld 2 (Schüler/in B)} \quad Name: \rule{4.5cm}{0.4pt}\\[0.15cm]
        \small\textbf{1.} Notiere hier für dich \textbf{5--8 zentrale Begriffe}.\\[0.05cm]
        \small\textbf{2.} Markiere in den anderen Feldern: \textbf{!} = tragend, \textbf{?} = unklar.
    };

    % --- FELD 3: UNTEN RECHTS ---
    \node[anchor=north west, text width=17.0cm, align=left, font=\sffamily] at (20.5, 6.4) {
        \textbf{\color{primaryBlue}\large Außenfeld 3 (Schüler/in C)} \quad Name: \rule{4.5cm}{0.4pt}\\[0.15cm]
        \small\textbf{1.} Notiere hier für dich \textbf{5--8 zentrale Begriffe}.\\[0.05cm]
        \small\textbf{2.} Markiere in den anderen Feldern: \textbf{!} = tragend, \textbf{?} = unklar.
    };

\end{tikzpicture}%
\end{document}
